\chapter{Исследовательская часть}

В данном разделе описана проверка корректности работы разработанного компилятора на тестовых программах.

\section{Цель исследования}
Цель исследования -- проверить, что компилятор корректно выполняет синтаксический и семантический анализ программ, написанных на текущей версии грамматики Plankalkul, а также генерирует корректный код CIL.

\section{Программа вычисления чисел Фибоначчи}
Первая программа реализует вычисление числа Фибоначчи.
Она демонстрирует определение пользовательского типа, работу с промежуточными переменными, \texttt{W}-цикл, guard-оператор и арифметические выражения.

\includelisting
{test_fibonacci.pk}
{Программа вычисления числа Фибоначчи}

Для проверки использованы команды:
\begin{verbatim}
uv run python -m plankalkulsharp parse
  tests\programs\test_fibonacci.pk
uv run python -m plankalkulsharp analyze
  tests\programs\test_fibonacci.pk
uv run python -m plankalkulsharp compile
  tests\programs\test_fibonacci.pk
  -o build\PlankalkulProgram.il
.\build\bin\Release\net10.0\PlankalkulProgram.exe 10
\end{verbatim}

Синтаксический анализ завершился построением AST.
Семантический анализ завершился без ошибок и предупреждений.
В отчете анализатора зарегистрирован тип \texttt{int} как битовая последовательность длиной 8 и программа \texttt{P1 fib} с одним параметром и одним результатом.
После команды \texttt{compile} был создан файл \texttt{build\textbackslash PlankalkulProgram.il} и собран исполняемый файл \texttt{build\textbackslash bin\textbackslash Release\textbackslash net10.0\textbackslash PlankalkulProgram.exe}.
Запуск полученной программы с аргументом \texttt{10} вывел значение \texttt{55}, что соответствует десятому числу Фибоначчи при нумерации, используемой в тестовой программе.

\section{Программа разворота связного списка}
Вторая программа демонстрирует более сложный набор возможностей языка.
В ней используются пользовательские типы \texttt{int}, \texttt{pointer}, \texttt{node}, \texttt{node\_pool}, массив записей, несколько программ, вызовы программ, \texttt{FOR}-цикл, \texttt{W}-цикл и компонентный доступ.

\includelisting
{test_reverse_list_fragment.pk}
{Фрагмент программы разворота связного списка}

Для проверки использованы команды:
\begin{verbatim}
uv run python -m plankalkulsharp parse
  tests\programs\test_reverse_list.pk
uv run python -m plankalkulsharp analyze
  tests\programs\test_reverse_list.pk
uv run python -m plankalkulsharp compile
  tests\programs\test_reverse_list.pk
  -o build\PlankalkulProgram.il
.\build\bin\Release\net10.0\PlankalkulProgram.exe 5 5
\end{verbatim}

Синтаксический анализ завершился построением AST.
Семантический анализ завершился без ошибок и предупреждений.
В отчете анализатора зарегистрированы записи, массив записей и программы \texttt{P0 test}, \texttt{P1 alloc\_pool}, \texttt{P4 create\_list}, \texttt{P5 reverse\_list}.
После компиляции была собрана та же исполняемая точка входа \texttt{PlankalkulProgram.exe}, но уже с CIL-кодом программы разворота списка.
Запуск с аргументами \texttt{5 5} вывел \texttt{4}: для списка длины 5 после разворота узел со значением 5 указывает на узел со значением 4, что проверяет изменение связей между элементами массива записей.

\section{Результаты проверки}
В таблице \ref{tbl:research_results} приведены результаты проверки актуальных программ.

\includelisting
{exe_test_output.txt}
{Фрагмент консольного вывода при компиляции и запуске актуальных программ}

\begin{table}[H]
	\caption{Результаты проверки актуальных программ}
	\label{tbl:research_results}
	\begin{tabular}{|p{4.8cm}|p{3.8cm}|p{4.7cm}|}
		\hline
		Программа & Проверяемые возможности & Результат \\ \hline
		\texttt{test\_fibonacci.pk} & Типы, присваивания, \texttt{W}-цикл, guard, арифметика & Разбор и анализ без ошибок; сборка exe успешна; запуск с аргументом \texttt{10} выводит \texttt{55} \\ \hline
		\begin{tabular}[t]{@{}l@{}}\texttt{test\_reverse}\\\texttt{\_list.pk}\end{tabular} & Записи, массивы, компонентный доступ, вызовы программ, \texttt{FOR}, \texttt{W} & Разбор и анализ без ошибок; сборка exe успешна; запуск с аргументами \texttt{5 5} выводит \texttt{4} \\ \hline
	\end{tabular}
\end{table}


\section*{Вывод по исследовательской части}
В результате исследования показано, что тестовые программы на языке Plankalkul успешно проходят синтаксический и семантический анализ, компилируются в CIL и собираются в исполняемый файл .NET.
Запуск полученных программ подтвердил работоспособность сгенерированного кода для арифметического цикла и для операций над массивом записей, моделирующим связный список.
